\section{Methods}
\label{sec:methods}

\subsection{Agent architecture}
\label{sec:architecture}

Figure~\ref{fig:architecture} shows the system. A runner turns a scenario into the task prompt, creates a fresh workspace, draws a sampling seed and starts a single tool-calling agent. The agent is built with LangChain's \texttt{create\_agent}, which compiles a two-node LangGraph loop: the model node calls the LLM, and the tools node executes the tool calls that the model returned, until the model answers without a tool call. This is the ReAct pattern~\citep{yao2023react} realised through native function calling: the JSON schemas of the level's tools are sent with every request, and the model-specific tool-call parser of vLLM converts the model output into structured calls, so the client never parses free text. We do not impose an explicit ``Thought:'' format, and thinking modes are switched off where a model has one; any reasoning appears in the ordinary text of the model's reply.

\begin{figure}[t]
  \centering
  \begin{tikzpicture}[
      font=\scriptsize,
      box/.style={draw, rounded corners=2pt, align=center, inner sep=3pt,
                  minimum height=10mm, text width=22mm},
      wide/.style={box, text width=28mm},
      flow/.style={-{Stealth[length=1.6mm]}, semithick},
      both/.style={{Stealth[length=1.6mm]}-{Stealth[length=1.6mm]}, semithick},
      aux/.style={-{Stealth[length=1.6mm]}, densely dashed},
    ]
    \node[box]  (scenario)  at (0, 0)       {\textbf{Scenario}\\topic, date window};
    \node[box]  (runner)    at (3.0, 0)     {\textbf{Runner}\\task prompt, seed,\\fresh workspace};
    \node[box]  (model)     at (6.3, 0.95)  {\textbf{Model node}\\one LLM call};
    \node[box]  (tools)     at (6.3, -0.95) {\textbf{Tools node}\\runs the tool calls};
    \node[wide] (vllm)      at (10.0, 0.95) {\textbf{vLLM server}\\OpenAI-compatible API,\\native tool calling};
    \node[wide] (toolset)   at (10.0, -0.95){\textbf{Tools of one level}\\L1: 12 \quad L2: 4 \quad L3: 1};
    \node[wide] (services)  at (10.0, -3.0) {\textbf{Services}\\arXiv API and cache,\\PDF text, workspace files};
    \node[wide] (review)    at (6.3, -3.0)  {\textbf{Review file}\\\texttt{reviews/review.md}};
    \node[box]  (evaluator) at (3.0, -3.0)  {\textbf{Evaluator}\\10 binary checks};
    \node[box]  (record)    at (0, -3.0)    {\textbf{Run record}\\JSONL and MLflow trace};
    \begin{scope}[on background layer]
      \node[draw, densely dashed, rounded corners=3pt, fill=black!4, inner sep=4pt,
            fit=(model)(tools)] (agent) {};
    \end{scope}
    \node[font=\scriptsize\bfseries, anchor=south] at (agent.north) {ReAct agent (LangGraph)};

    \draw[flow] (scenario) -- (runner);
    \draw[flow] (runner) -- (runner -| agent.west) node[midway, above] {task};
    \draw[flow] ([xshift=-4mm]model.south) -- ([xshift=-4mm]tools.north)
                node[midway, left] {calls};
    \draw[flow] ([xshift=4mm]tools.north) -- ([xshift=4mm]model.south)
                node[midway, right] {results};
    \draw[both] (model) -- (vllm);
    \draw[flow] (tools) -- (toolset);
    \draw[flow] (toolset) -- (services);
    \draw[aux]  ([xshift=11mm]toolset.north) -- ([xshift=11mm]vllm.south)
                node[midway, left, align=right] {LLM calls\\inside tools};
    \draw[flow] (services) -- (review);
    \draw[flow] (review) -- (evaluator);
    \draw[flow] (evaluator) -- (record);
    \draw[aux]  (runner.south) -- ++(0, -0.45) -| (record.north)
                node[pos=0.25, above] {tokens, calls, time};
  \end{tikzpicture}
  \caption{System architecture. The runner turns a scenario into the task prompt and starts the agent, a LangGraph loop that alternates between model calls served by vLLM and the execution of tool calls. Only the set of tools changes between levels; LLM-backed tools call the same model (dashed). After the agent stops, the evaluator scores the review file, and the run record stores the checks, the cost counters collected during the run (dashed) and the configuration.}
  \label{fig:architecture}
\end{figure}

A run ends in one of five states. It is \emph{completed} when the agent stops on its own. It hits the \emph{step limit} after 100 model calls of the agent loop (LLM calls made inside tools do not count). It is stopped as a \emph{loop} when the agent requests the same tool call with the same arguments for the eighth time in a row, before that call runs, or when ten of its tool calls in a row have failed, before the model is called again; such runs could only have ended at the step limit, and stopping them early saves compute but also lowers their token counts. It has \emph{failed} when the context window overflows or when the final reply is cut off at the 8{,}192-token completion limit without a tool call. It ends with an infrastructure \emph{error} when, for example, arXiv still cannot be reached after the scoring of the review has been retried three times (after 1, 3 and 10 minutes; the agent run itself is never repeated for this); such runs are excluded and repeated. A failing tool call is not an error in this sense: its error message is returned to the model, which can correct its input or choose another approach, as the system prompt instructs.

\subsection{Tool granularity levels}
\label{sec:levels}

The three levels cover the same workflow, partitioned into tools of different size (Table~\ref{tab:tools}). All tools are thin wrappers over one service layer (arXiv client, PDF pipeline, LLM operations and workspace files). Besides how much of the workflow a single call performs, the levels differ in four respects that are part of the design: Level~1 offers four distractor tools and a tool to read workspace files, which Levels~2 and~3 lack; Level~2 search results carry ready APA references but no abstracts; at Level~1 the agent decides how much of each paper to read, whereas Levels~2 and~3 always summarise the first six chunks; and Level~3 decides the papers, their number and the file structure in code.

\begin{table}[t]
  \caption{Tools per granularity level. Tools marked $^\ast$ call the LLM internally.}
  \label{tab:tools}
  \centering
  \small
  \begin{tabularx}{\linewidth}{@{}l>{\raggedright\arraybackslash}X@{}}
    \toprule
    Level & Tools \\
    \midrule
    1 (12 micro-tools) & \texttt{search\_arxiv}, \texttt{search\_wikipedia}, \texttt{download\_pdf}, \texttt{extract\_text\_from\_pdf}, \texttt{split\_text\_into\_chunks}, \texttt{summarize\_chunk}$^\ast$, \texttt{extract\_keywords}$^\ast$, \texttt{translate\_text}$^\ast$, \texttt{format\_apa\_citation}, \texttt{create\_markdown\_table}, \texttt{write\_file}, \texttt{read\_file} \\
    \addlinespace
    2 (4 medium tools) & \texttt{search\_papers} (arXiv search that returns APA references), \texttt{get\_and\_summarize\_paper}$^\ast$ (download, extract, split and summarise one paper), \texttt{write\_review} (write the whole review file), \texttt{append\_bibliography} (write the bibliography for given identifiers) \\
    \addlinespace
    3 (1 macro-tool) & \texttt{conduct\_literature\_review}$^\ast$ (search, read, summarise and write the complete review from a single query) \\
    \bottomrule
  \end{tabularx}
\end{table}

\paragraph{Level 1 (micro-tools).} Every tool performs one small operation, so the agent must plan the whole chain: search, download each PDF, extract its text, split it, summarise the chunks, format the references and write the review with \texttt{write\_file}, which replaces the file's content. Four tools are not needed for the base task (Wikipedia search, keyword extraction, translation and Markdown tables); they act as distractors there and become useful in the variants with an added requirement. How much of a paper the agent reads is its own choice: \texttt{summarize\_chunk} accepts at most one chunk of 3{,}000 characters and asks the agent to split longer texts first, and every text that a tool shortens ends with a marker giving the characters shown and the total.

\paragraph{Level 2 (medium tools).} One call of \texttt{get\_and\_summarize\_paper} downloads a paper, extracts and splits its text, summarises up to six chunks and combines the chunk summaries into one summary of a few sentences (up to seven internal LLM calls); \texttt{write\_review} replaces the review file with the agent's text; and \texttt{append\_bibliography} writes an alphabetical APA bibliography after checking that every identifier exists on arXiv. The agent still chooses the papers and composes the review.

\paragraph{Level 3 (macro-tool).} The single tool takes only a query. It retrieves $2N=10$ candidates, summarises them in order of relevance with the Level~2 pipeline, skipping papers whose PDF cannot be read, until $N$ papers are summarised (up to 35 internal LLM calls), and writes the complete review with its bibliography. Its result states that the review is complete, or how many papers are missing, and each call starts from a clean workspace. The workflow is fixed in code; the agent decides only the query.

The levels also change the prompt: the JSON schemas of the tools take 4{,}826, 1{,}750 and 621 characters at Levels~1, 2 and~3. This prompt overhead is confounded with the level (Section~\ref{sec:variables}).

\subsection{Models}
\label{sec:models}

We evaluate five open-weight, instruction-tuned models (Table~\ref{tab:models}) in the checkpoints listed there; the three larger ones are 4-bit checkpoints, Qwen3.5-4B and Qwen3.5-0.8B run unquantised in BF16. They form two comparisons that share the reference model Qwen3.5-9B, which is also the baseline model: a \emph{family} comparison of three models of similar size from different families, and a \emph{scale} comparison of three sizes within the Qwen3.5 family. Each model uses its own vLLM tool-call parser, and Gemma and Llama use dedicated tool chat templates. Every model is limited to one tool call per turn, because the Llama template cannot render more and one policy keeps the number of calls comparable. All models share explicit sampling settings, sent with every request so that no model falls back on the defaults of its own generation config: temperature 0.7, top-$p$ 0.95, top-$k$ 20, a random seed per run that is recorded, and at most 8{,}192 completion tokens. Parameter counts are nominal. Within the scale series the 9B model is 4-bit while the two smaller ones are unquantised, and within the family series Gemma uses a quantisation-aware-trained checkpoint while Qwen and Llama use post-training AWQ; we treat both as limitations. The served checkpoint, quantisation, parser, chat template and vLLM version are part of every run's configuration hash.

\begin{table}[t]
  \caption{Models. All are served by vLLM~0.33.0 with a 32{,}768-token context window.}
  \label{tab:models}
  \centering
  \small
  \begin{tabular}{@{}llrll@{}}
    \toprule
    Model & Comparison & Params & Quantisation & Tool-call parser \\
    \midrule
    Qwen3.5-9B (baseline) & family, scale & 9B & AWQ 4-bit & \texttt{qwen3\_coder} \\
    Gemma 4 12B IT & family & 12B & QAT W4A16 & \texttt{gemma4} \\
    Llama 3.1 8B Instruct & family & 8B & AWQ INT4 & \texttt{llama3\_json} \\
    Qwen3.5-4B & scale & 4B & none (BF16) & \texttt{qwen3\_coder} \\
    Qwen3.5-0.8B & scale & 0.8B & none (BF16) & \texttt{qwen3\_coder} \\
    \bottomrule
  \end{tabular}
\end{table}

\subsection{Memory, planning, verification and communication}
\label{sec:mechanisms}

\paragraph{Memory.} The agent has only short-term memory: the full message history of the run, without truncation or summarisation, bounded by the context window; nothing carries over between runs. The workspace serves as external working memory: tools exchange relative file paths (for example \texttt{papers/<id>.pdf} or \texttt{summaries/<id>.md}) instead of copying document text into the conversation, and reading a file returns at most 6{,}000 characters, followed by a marker when the file is longer.

\paragraph{Planning.} There is no separate planner. The plan emerges from the ReAct loop, guided by the tool descriptions and by the task prompt, which is the same at every level and lists only the properties of the finished review (number of papers, date window, summaries that reflect each paper's full text, file structure), not the steps to produce it. At Level~2 part of the plan (fetch, extract, split, summarise) sits inside one tool, and at Level~3 the whole plan is fixed in the macro-tool.

\paragraph{Verification.} Tool arguments are validated against schemas derived from the tool signatures, including value ranges (for example 1--10 search results), and lists passed as JSON strings are decoded. Tools check arXiv identifiers, the workspace sandbox, file types and sizes and the PDF format, and they report failures as messages the agent can act on; \texttt{append\_bibliography} also verifies that each identifier exists. There is no LLM-based self-critique, and the agent never sees the evaluator's verdict, which is computed only after the run ends.

\paragraph{Communication.} The system has a single agent. LLM-backed tools call the same model on the same server with the same sampling settings and seed; their token usage is attributed to the tool and included in the cost of the run (Section~\ref{sec:cost}).

\subsection{Outcome evaluation}
\label{sec:evaluation}

After the agent stops, \texttt{reviews/review.md} is scored in three layers, because a single set of format checks proved both too strict and too blind (Section~\ref{sec:protocol-changes}): at Level~3 the code writes the format, so every Level~3 review passed it, including one on an unrelated topic, while reviews of five correct, read papers failed nine of ten checks over a heading level.

\paragraph{Format checks.} The ten binary checks of Table~\ref{tab:checks} verify the required structure and references; they parse paper sections from second-level headings and expect exact APA lines. We report them as \emph{format compliance}.

\paragraph{Fact checks.} Four checks take every arXiv identifier anywhere in the file, whatever its formatting, and verify facts that need no interpretation: at least $N$ distinct papers are cited, every identifier exists, every paper was submitted within the window, and the text of every cited paper was extracted in the run's own workspace. Facts about papers are looked up on arXiv through the cache.

\paragraph{Requirement checks.} For a variant, the \texttt{table} requirement passes when a Markdown table between the title and the first paper section has a row for every cited paper (by title or identifier), and the \texttt{polish} requirement when the summary text of every paper section, without titles, references and URLs, is detected as Polish (\texttt{langdetect}, fixed seed). The \texttt{no\_surveys} requirement is decided by the judge.

\paragraph{LLM judge.} After all runs, Qwen3.5-9B, served like the evaluated models, rates every scored review on three 1--5 scales with short anchored rubrics: whether it is a valid review with one section and summary per paper, whether the papers fit the topic, and the faithfulness of two summaries, chosen with a seed derived from the run, against the first 12{,}000 characters of the paper's extracted text. For \texttt{no\_surveys} variants it also decides, from each cited paper's arXiv title and abstract, whether the paper is a survey. The judge sees neither the model, the level nor the run; it samples with temperature 0 and a fixed seed, and its answers are constrained by guided decoding to schemas that contain only integers or booleans, without justifications; an answer that does not fit is retried once and otherwise recorded as missing. The judge is also one of the evaluated models, so it may favour its own style; blinding reduces this risk but cannot remove it, and the judgments were not validated against human ratings.

\begin{table}[t]
  \caption{Binary checks applied to the review file. If the file is missing, all checks fail.}
  \label{tab:checks}
  \centering
  \small
  \begin{tabularx}{\linewidth}{@{}l>{\raggedright\arraybackslash}X@{}}
    \toprule
    Check & Passes when \\
    \midrule
    \texttt{file\_exists} & \texttt{reviews/review.md} exists \\
    \texttt{section\_count} & there are exactly $N$ paper sections (second-level headings other than the bibliography) \\
    \texttt{summaries\_present} & every paper section contains at least 40 characters of text besides its reference \\
    \texttt{citations\_apa} & every reference line matches the APA~7 pattern for arXiv preprints \\
    \texttt{ids\_exist} & every cited arXiv identifier resolves to a paper \\
    \texttt{dates\_in\_window} & the first version of every paper was submitted within the scenario window \\
    \texttt{titles\_match} & the real title of every paper appears in its heading and in its reference \\
    \texttt{bibliography\_matches} & the bibliography lists exactly the cited identifiers \\
    \texttt{no\_duplicates} & no paper is cited twice \\
    \texttt{papers\_read} & the run extracted the text of every cited paper (how much of it was read is measured separately) \\
    \bottomrule
  \end{tabularx}
\end{table}

\paragraph{Outcomes.} Each run has two outcome axes. Its review is \emph{artifact valid} when every fact check passes, every judge rating is at least 4 and a variant meets its requirement, whatever the run's status. The run is \emph{reliable} when the agent ended it itself (it neither hit the step limit, nor was stopped as a loop, nor failed). A \emph{success} is both. The original rule---completed and all ten format checks passed---is kept as \emph{format success}, and the \emph{score} is the fraction of passed format checks. Runs with an infrastructure error are excluded from all metrics and repeated.

\subsection{Cost measurement}
\label{sec:cost}

LangChain callbacks meter every run. For each LLM call they record the input and output tokens reported by the server and attribute the call to the agent loop or, if it was issued inside a tool, to that tool; they also count tool calls and tool errors. Tool-internal calls are essential for a fair comparison: at Level~3 almost all LLM work happens inside the macro-tool, so counting only the agent's own calls would make coarse tools look cheap. For the main experiment the meter also records (i) the time spent in LLM calls, separately for the agent and for tools, and the remaining tool time, which is dominated by network access; (ii) the input tokens of the first agent call, which measure the prompt overhead of a level; (iii) malformed tool calls, such as unparseable arguments or unknown tool names, which never start a tool run, and replies that carry a tool call as plain text because the server's parser missed it; (iv) tool errors split into environment errors (network, arXiv HTTP errors, missing PDFs) and agent errors; (v) per paper, the number of distinct characters of its extracted text that reached the agent or a model; and (vi) the arXiv requests and PDF downloads that missed the shared caches. Wall-clock time also includes arXiv throttling, so we report it only as an indication.

\subsection{Protocol changes after the pilot}
\label{sec:protocol-changes}

A methodology review of a 150-run diagnostic screen (one repeat of every model, level and base scenario) led to the following changes before the main experiment; the screen's results are not reported. (1)~Sampling: top-$p$ and top-$k$ are now sent explicitly and identically to every model, where previously each model used its own defaults (the Qwen models without any truncation). (2)~One tool call per turn for every model, where previously only Llama was limited. (3)~Qwen3.5-4B is served in BF16 instead of online FP8. (4)~The task prompt states the outcome instead of prescribing tool steps; the old wording made the Level~3 agent re-run its finished macro-tool. (5)~Shortened text is marked and \texttt{summarize\_chunk} rejects inputs longer than one chunk, where previously Level~1 silently read only the first 6{,}000 characters of a paper. (6)~Levels~2 and~3 combine chunk summaries into one summary instead of concatenating them. (7)~Level~1 overwrites files instead of only appending, which had duplicated sections. (8)~Tool descriptions no longer show a real-looking arXiv identifier as an example, which weak models had copied. (9)~Caches are filled before the run, missing PDFs are cached and tried at their latest version, scoring is retried on arXiv outages, and the configuration hash covers the served model. (10)~The loop guard stops at the eighth identical call instead of the fifth. Every run is traced in MLflow with its parameters, metrics, review file and per-check assessments.
